\documentclass[11pt]{article}
\usepackage[margin=1in]{geometry}
\usepackage{amsmath,amssymb,amsthm,microtype,hyperref}
\hypersetup{hidelinks,pdfauthor={T},pdftitle={Quorum Intersection and Fenced Cutover for Agent Tools}}
\newtheorem{theorem}{Theorem}
\title{Quorum Intersection and Fenced Cutover for Agent Tools}
\author{T}
\date{}
\begin{document}
\maketitle
\begin{abstract}
Two configurations can each satisfy their own quorum rule while
permitting disjoint certificates across a membership change.
We give exact minimum intersections for direct, staged, and
joint threshold certificates. A single median determines the
joint minimum. For two stages that may coexist, these minima
determine whether the common signers must include someone outside
the fault budget. Exhaustive signer-set checks verify the formulas
on small committees. A separate finite model tests how a sink
can end old authority before accepting a new epoch while
preserving the record of earlier irreversible effects. For
collective effects, closure and prior absence depend on different
families of sink sets. Exact evidence cuts identify when partial
acknowledgments suffice. A local SQLite experiment exercises
transaction ordering and interrupted acknowledgments.
\end{abstract}

\section*{Summary and scope}
\paragraph{Question.} How much overlap and terminal evidence are needed
when changing the signers that authorize an effect?
\paragraph{Finding.} Exact threshold intersections characterize whether
certificate pairs must share an honest signer. Sink closure prevents
later old effects. Final occupancy evidence is separately needed to avoid
duplicating an earlier irreversible effect.
\paragraph{Evidence and scope.} The overlap formulas and cutover
rule are supported by exact set and interleaving checks. They apply to
authorization gates that enforce durable signer and sink rules, without
establishing a deployed consensus protocol.

\section{What needs to intersect}
Let the old committee $A$ have size $m$ and the new committee
$B$ have size $n$, with $s=|A\cap B|$. The old and new thresholds
are $a$ and $b$, respectively, each positive and no greater than
its committee size. A joint certificate is a subset of $A\cup B$
with at least $a$ signers in $A$ and at least $b$ in $B$. A signer
in the intersection contributes to both thresholds.

Raft uses majorities in both configurations during joint consensus
\cite{raft}. Byzantine quorum systems require an intersection
containing sufficiently many correct signers \cite{malkhi}.
Here the same conflict domain is enforced across configurations,
at most $f$ signer identities may equivocate across that domain,
and every other signer durably refuses to sign conflicting
effects. Under these assumptions, a pair of certificate families
rules out conflicting certificates exactly when its minimum
intersection exceeds $f$. If the minimum is at most $f$, the
common signers in a minimizing pair may all equivocate and the
disjoint remaining signers can each sign one effect.

\section{Exact thresholds}
Write $[x]_+=\max(0,x)$, $p=m-s$, and $q=n-s$. Let $I_{AB}$,
$I_{AJ}$, $I_{JB}$, and $I_{JJ}$ denote the minimum signer
intersections for old/new, old/joint, joint/new, and joint/joint.

\begin{theorem}
The first three minimum intersections are
\begin{align*}
 I_{AB}&=[a+b-(m+n-s)]_+,\\
 I_{AJ}&=[a+\max(a,b-q)-m]_+,\\
 I_{JB}&=[b+\max(b,a-p)-n]_+.
\end{align*}
For the joint/joint minimum define $h=\max(0,a-p,b-q)$,
$\alpha=2a-p$, and $\beta=2b-q$. Let $T_*$ be the median
of $(s,\alpha,\beta)$ clamped to $[2h,2s]$. Then
\[
 I_{JJ}=[T_*-s]_++[\alpha-T_*]_++[\beta-T_*]_+ .
\]
\end{theorem}
\begin{proof}
For direct overlap, fill the exclusive parts first. The smallest
necessary shared counts are $[a-p]_+$ and $[b-q]_+$. Two subsets
of the shared pool with those sizes can overlap by as little as
their summed size minus $s$, floored at zero. This simplifies to
the formula for $I_{AB}$.

A joint certificate must contain at least $\max(a,b-q)$ signers
from $A$. Fill $B$'s exclusive part as far as possible, then meet
the remaining $B$ threshold from the shared pool and complete the
$A$ threshold. This attains the lower bound. An old certificate
can take as many of its signers as possible from the complement
of that set in $A$, which proves the formula for $I_{AJ}$.
Exchanging the committee roles gives $I_{JB}$.

For two joint certificates with shared counts $c_1,c_2$, the
exclusive counts can be chosen as $[a-c_i]_+$ in the old
pool and $[b-c_i]_+$ in the new pool. The shared counts range
from $h$ to $s$. Let $T=c_1+c_2$. The least overlap in the
shared pool is $[T-s]_+$. The sum of exclusive old counts
is at least $[2a-T]_+$, and the sum of exclusive new counts
is at least $[2b-T]_+$. Both bounds are attained together by
taking $c_1=\lfloor T/2\rfloor$ and $c_2=\lceil T/2\rceil$.
For each pool we can choose the two subsets to attain their
minimum overlap, independently of the other pools. It follows
that the exact minimum for fixed $T$ is
$[T-s]_++[\alpha-T]_++[\beta-T]_+$.
The slope starts at $-2$ and increases by one at each
breakpoint $s,\alpha,\beta$, counting multiplicities. A median
minimizes the expression, and clamping it to the feasible
interval preserves optimality. All endpoints and breakpoints
are integers, so the minimum is attained at a feasible integer.
\end{proof}

\section{Availability and a single equivocator}
A joint certificate remains possible under every withholding of
at most $f$ signers, with $f<\min(m,n)$, exactly when
$a\le m-f$ and $b\le n-f$. Sufficiency follows because enough
signers remain in each committee. Necessity follows by placing
all $f$ omissions inside one committee. Among thresholds meeting
these inequalities, the largest ones maximize the direct
minimum, giving $[s-2f]_+$. Thus concurrent old and new
authorization with a guaranteed honest common signer under
these availability conditions requires $s\ge 3f+1$.

The same-epoch minima are $[2a-m]_+$ for two old certificates
and $[2b-n]_+$ for two new certificates. Suppose a staged
cutover permits old and joint certificates to coexist, then
joint and new, and a fence rules out old and new coexistence.
With the largest available thresholds $a=m-f$ and $b=n-f$,
the old/joint and joint/new minima equal the two same-epoch
minima. The joint/joint minimum is at least their maximum.
Hence $m,n\ge3f+1$ is necessary and sufficient for all allowed
pairs to have an honest intersection under the availability
requirement, regardless of shared membership. The necessity
comes from each configuration's own pair of certificates.
The conclusion relies on enforcing the permitted coexistence
schedule at the sinks.

For the old committee $\{A,B,C\}$ and the new committee
$\{C,D,E\}$ with majority thresholds, the exact minima
in the order above are $(0,1,1,1)$. The two valid joint
certificates $\{A,C,D\}$ and $\{B,C,E\}$ intersect only
at $C$. A single equivocator can therefore defeat the
honest-intersection requirement even for two joint certificates.
For disjoint four-member committees with threshold three,
the minima are $(0,2,2,4)$. With one potentially faulty
signer, adjacent phases have an honest intersection, but old
and new alone still have none.

\section{Closing the old epoch at the effect sinks}
Chubby's sequencers attach a generation to a lock and allow
the effect server to reject stale use \cite{burrows}. In the
following model, the authorization service knows the complete
set of sinks that may still commit an old effect. It freezes
issuance of old rights or enforces a ceiling on every old
right's lifetime before activation. Each trusted sink checks
the epoch and expiry at its irreversible commit point.
The adversary may delay requests, evidence, and epoch messages.

For sink $i$, let $e_i$ be a certified hard upper cutoff,
so all its old commits occur strictly before $e_i$. Let $b_i$
be the time at which the coordinator first receives proof of
the sink's irreversible old-epoch fence. The fence and commit
decisions share a durable order. A missing bound is represented
by $+\infty$. Under these evidence rules, the earliest time
at which all old authority is known closed is
\[
 t_{\rm close}=\max_i \min(e_i,b_i).
\]
At that time each sink has either passed its cutoff or supplied
proof of a fence. Before that time there is a sink with
neither form of evidence. If an otherwise eligible old request
can still commit there, delaying the fence and delivering the
request yields a history consistent with what the coordinator
knows. This makes the bound exact for these observations when
no other mechanism can exclude such a request. When the
coordinator knows only an interval for the current time, it
uses the lower endpoint to establish that expiry has passed.

For example, take two sinks with hard old cutoffs $(9,13)$
and received fence proofs at times $(11,8)$. The coordinator
has closure evidence for both at time nine. A prior absence
report from the first sink may still be stale, since it could
have committed between the report and time nine. This remains
possible even if the old and new signer sets each met their
own quorum rule.

Closing the old epoch prevents future old commits. A new
commit that must be unique across both epochs also needs
evidence about the past. For each sink, require a terminal
record for the conflict key. Such a record reports whether
an old effect already committed and is produced atomically
with the fence or from an authoritative snapshot after
closure. Bind it to the sink, epoch, and conflict key.
Only a full set of absent records admits the new effect.
An unknown record leaves the key blocked, even after expiry.

To see sufficiency, an absent terminal record rules out a
prior commit at its sink, and closure rules out any later old
one. The argument applies to every sink that could have
committed the key. For necessity in this observation model,
a missing terminal record can leave two indistinguishable
histories, one with a prior commit and one without it. An
absence observation from before closure may be followed by a
commit before the fence or expiry. The same test can use an
authoritative shared journal if it provides equivalent
complete evidence. The rule concerns uniqueness of the
specific irreversible effect. It does not undo a committed
effect or grant new authority on its own.

\subsection{Partial evidence for collective effects}
Some tools can commit an effect only when several sinks remain
live under the old epoch. Let $\mathcal H$ be a nonempty family
of nonempty sink sets. Every old commit requires all the members
of some $H\in\mathcal H$ to be live at its irreversible commit
point. Write $c_i=\min(e_i,b_i)$ for the time when closure at sink
$i$ is certified. A missing certificate and cutoff give $c_i=+\infty$.

Past effects need a separate assumption. Let $\mathcal R$ be the
nonempty family of possible durable record footprints for one old
commit on the conflict key. Each footprint is nonempty. Write $a_i$
for the time a trustworthy terminal absence record from sink $i$
is available. An unknown or positive record has $a_i=+\infty$.
The absence record covers the key and epoch and rules out a durable
old commit mark at that sink. It must cover the entire old epoch,
including in-flight effects and any delayed replication of the mark.

\begin{theorem}
Under these evidence rules, closure and prior absence are certified
respectively by times
\[
 \max_{H\in\mathcal H}\min_{i\in H}c_i,
 \qquad
 \max_{R\in\mathcal R}\min_{i\in R}a_i.
\]
The first bound is exact when each unhit live set can still
commit in some history consistent with the observations. The second
is exact when a footprint disjoint from all observed terminal absences
can still belong to a past commit consistent with those observations.
\end{theorem}
\begin{proof}
After the first bound, the closed sinks meet every required live set.
An old commit would require a live sink in that closed set, which is
impossible. Before the bound, an unhit live set remains. The stated
realizability assumption supplies a continuation with an old commit.
For past effects, every possible footprint meets the set of observed
terminal absences after the second bound. A prior commit would then
contradict an absence record. Before that bound, there is a footprint
containing no such record, and the realizability assumption supplies
a compatible old commit. The first time every set is hit is the latest
of the first-hit times, giving both formulas.
\end{proof}

With $n$ sinks and every $q$-sink subset allowed as a live set,
$n-q+1$ closure proofs suffice. In a five-sink example, take
$q=3$ and closure times $(2,5,8,11,+\infty)$. Closure is certified
at time eight. If every old effect leaves a durable mark on some
two-sink subset, and terminal absence times are $(3,6,9,12,+\infty)$,
prior absence is certified at time twelve. These tests together
permit a unique new effect at time twelve if the new authority
and its other gates already hold. A sink can remain unreachable.

Cached signatures cannot count as live participation at commit.
For a two-out-of-three gate, a sink can sign at time one and close
at time two. Another closes at time three. At time four the last
sink can add a valid old signature to the saved one. A gate that
checks only the signatures can still accept the old effect.
The live-set calculation applies only if the commit protocol
itself enforces the required epoch checks at linearization.

\section{Reproducible checks}
The original three-by-three fixture had five disjoint pairs
among nine old/new majority pairs. The generalized check
compared the closed forms with enumerated signer sets on 74
committee layouts. All 3,216 minimum-overlap comparisons and
3,468 availability checks agreed. This is a finite check of
the certificate arithmetic.

A separate state search modeled two sinks, one conflict key,
finite deadlines, and delayed terminal records. The proposed
rule reached 2,677 states with neither an old commit after
activation nor a new commit duplicating an old one. Requiring
closure from only one sink and allowing a fence to roll back
each admitted an old effect after activation. Taking expiry
as proof of prior absence and using pre-closure absence each
admitted a duplicate. The receipt contains a shortest witness
for each change to the model. The search does not cover
network liveness or the correctness of a physical effect sink.

The collective-effect check compared both time bounds with a
stepwise oracle on 512 cases. A separate enumeration of possible
record footprints checked 6,336 absence patterns. It certified
2,899 patterns with fewer than all sinks reporting. A cached-vote
control passed the signature threshold after live closure, as
predicted. These counts describe finite models of the declared
commit and record rules.

\subsection{A local transaction and crash experiment}
We implemented the effect boundary as a database row with a
conflict key. An effect starts a SQLite write transaction,
checks the epoch and key, and inserts the row before committing.
Fencing updates the accepted epoch and reads the key in another
write transaction. The coordinator receives the terminal record
only after that transaction commits. In the test, the committed
row is the complete effect.

We ran the same cases with SQLite 3.54.0 in rollback-journal
and write-ahead-log modes, with full synchronous writes. An old
transaction that acquired the write lock first prevented the
fence from acquiring it. After the old commit, the fence observed
the occupied key. Once the fence committed, the old epoch could
no longer insert a different key. A fresh authorized key could
still commit under the new epoch.

Eight child processes exited at specified points before or after
an effect or fence commit. Within each pair, the caller saw the
same response and exit code. Reopening the store distinguished
the outcomes. A pre-commit effect exit left the key absent.
A post-commit exit left it occupied. Likewise, reopening after
a post-commit fence exit preserved the new epoch and rejected
an old request. Reopening after a pre-commit fence exit retained
the old epoch until a later successful fence.

Two negative controls exercised the modeled failures. A process
that cached authority before fencing could insert an old-epoch row
after the fence if its later transaction skipped the epoch check.
A coordinator that inferred absence from a lost acknowledgment
could commit the same key in a separate store, whose local key
index could not see the original row. The correct route admitted
a new write only after a terminal absence record from the old
store. These tests cover process exit and local database state.
They leave authentication, hardware power loss and effects outside
the database to separate experiments.

\section{Application and next test}
An agent tool that authorizes sensitive writes can compare the
minimum overlap of every certificate pair its sink may accept.
If the tool uses staged authorization with small shared membership,
it needs a durable fence that prevents old-only certificates
from authorizing effects once new-only authority is active.
Replayed certificates, disconnected sinks, and signer-state
loss can invalidate that ordering. The arithmetic assumes the
same durable conflict rule at every phase and says nothing
about whether such a fence, authenticated key change, or
network progress guarantee has been implemented.

\paragraph{Review status.} Independent manuscript review is pending.

\begin{thebibliography}{9}
\bibitem{raft} Diego Ongaro and John Ousterhout. \emph{In Search of an Understandable Consensus Algorithm}. 2014. \url{https://raft.github.io/raft.pdf}.
\bibitem{malkhi} Dahlia Malkhi and Michael K. Reiter.
Byzantine Quorum Systems. \emph{Distributed Computing}
11, 203--213, 1998.
\url{https://users.ece.cmu.edu/~reiter/papers/1998/DC.pdf}.
\bibitem{burrows} Mike Burrows. The Chubby Lock Service for
Loosely-Coupled Distributed Systems. \emph{OSDI},
335--350, 2006.
\url{https://storage.googleapis.com/gweb-research2023-media/pubtools/4444.pdf}.
\end{thebibliography}
\end{document}
